\section{Grafiken}
\begin{frame}
\frametitle{Grafiken}
\framesubtitle{Einbinden von Grafiken} 
\begin{exampleblock}{Neue Pakete in diesem Abschnitt}
\begin{itemize}
\item graphicx 
\item float
\item cleveref
\end{itemize}
\end{exampleblock}

\begin{block}{Neue Befehle in diesem Abschnitt}
\begin{itemize}
\begin{ttfamily}
\item \color{nounibaredI}\textbackslash includegraphics\color{black}\{Datei\}
\item \color{nounibaredI}\textbackslash caption\color{black}\{Bildunterschrift\}
\item \color{nounibaredI}\textbackslash label\color{black}\{Schl"ussel\}
\item \color{nounibaredI}\textbackslash ref\color{black}\{Schl"ussel\}
\end{ttfamily}
\end{itemize}
\end{block}

\end{frame}

%-------------------------------------------------------------------------------

\begin{frame}
\frametitle{Grafiken}
\framesubtitle{Abbildungen einfügen}
\begin{tabbing}
\textbackslash includegraphics[option]\{Datei\}x\=\kill
\hspace{-10mm}
\begin{ttfamily}
\color{unibablueI}\textbackslash begin\color{black}\{figure\}\color{nounibagreenI}[option]\color{black}
\end{ttfamily}
\>\textbf{Mögliche Optionen zur Positionierung:}\\
\>\begin{ttfamily}\color{nounibagreenI}[h]\color{black}\end{ttfamily} = Hier an dieser Stelle\\
\>\begin{ttfamily}\color{nounibagreenI}[t]\color{black}\end{ttfamily} = Oben auf der Seite\\
\>\begin{ttfamily}\color{nounibagreenI}[b]\color{black}\end{ttfamily} = Unten auf der Seite\\
\>\begin{ttfamily}\color{nounibagreenI}[p]\color{black}\end{ttfamily} = Platzierung auf einer eigenen
Seite\\[5mm]
~\\[5mm]
\hspace{-10mm}
\begin{ttfamily}
\small\color{nounibaredI}\textbackslash includegraphics\color{nounibagreenI}[option]\color{black}\{datei\}
\end{ttfamily}
\>\textbf{Mögliche Optionen für das Einfügen:}\\
\>\begin{ttfamily}\color{nounibagreenI}[width=300pt]\color{black}\end{ttfamily} = Skalieren auf eine Breite\\
\>\begin{ttfamily}\color{nounibagreenI}[height=5cm]\color{black}\end{ttfamily} = Skalieren auf eine Höhe\\
\>scale, angle und noch viele mehr\ldots\\
\>Kombinationen möglich:\\
\>\begin{ttfamily}\color{nounibagreenI}[width=\textbackslash textwidth,height=5cm]\color{black}\end{ttfamily}
\end{tabbing}
\end{frame}

%-------------------------------------------------------------------------------

\begin{frame}
\frametitle{Grafiken}
\framesubtitle{Positionierung von Abbildungen}
\begin{columns}
\begin{column}{.5\textwidth}
{\ttfamily {\footnotesize
\color{nounibaredI}\color{nounibaredI}\textbackslash documentclass\color{black}\{article\} \\
\color{nounibaredI}\color{nounibaredI}\textbackslash usepackage\color{black}\{graphicx\} \\
\color{nounibaredI}\color{unibablueI}\textbackslash\color{unibablueI}begin\color{black}\color{black}\{document\} \\
\color{nounibaredI}\color{unibablueI}\textbackslash\color{unibablueI}begin\color{black}\color{black}\{figure\}\color{nounibagreenI}[h]\color{black} \\
\color{nounibaredI}\color{unibablueI}\textbackslash\color{unibablueI}begin\color{black}\color{black}\{center\} \\
	\color{nounibaredI}\color{nounibaredI}\textbackslash includegraphics\color{black}\color{nounibagreenI}[width=50mm]\color{black}\{tux.png\} \\
\color{nounibaredI}\color{nounibaredI}\textbackslash caption\color{black}\{Der kleine Tux\} \\
\color{nounibaredI}\color{nounibaredI}\textbackslash label\color{black}\{img:kleinertux\} \\
\color{nounibaredI}\color{unibablueI}\textbackslash\color{unibablueI}end\color{black}\color{black}\{center\} \\
\color{nounibaredI}\color{unibablueI}\textbackslash\color{unibablueI}end\color{black}\color{black}\{figure\} \\
\color{nounibaredI}\color{unibablueI}\textbackslash\color{unibablueI}end\color{black}\color{black}\{document\} \\
}}

\begin{alertblock}{Obacht!}
\texttt{\color{nounibaredI}\textbackslash label\color{black}\{\}} immer nach \texttt{\color{nounibaredI}\textbackslash caption\color{black}\{\}}
\end{alertblock}
\end{column}

\begin{column}{.5\textwidth} 
\begin{figure}
\begin{center}
    \includegraphics[width=\textwidth]{image/tux.png}
\caption{Der kleine Tux}
\label{img:kleinertux}
\end{center}
\end{figure}
\end{column}
\end{columns}
\end{frame}

%-------------------------------------------------------------------------------

\begin{frame}[t]
\medskip
\frametitle{Grafiken}
\framesubtitle{Positionierung von Abbildungen II }
\begin{block}{Exakte Positionierung erzwingen}
Package {\ttfamily \textbf{float}} und \textbf{[H]} verwenden\\
(ohne R"ucksicht auf entstehende Leerfl"achen)
\end{block}
%Trotz Definition einer Positionsumgebung verrutscht das Bild oft, da es nicht immer m\"oglich ist das Bild an passender Stelle einzufügen.\\
%\textbf{L\"osung:} Das Package {\ttfamily float} erzwingt die Positionierung an genau dieser Stelle im Text (allerdings ohne R"ucksicht auf entstehende Leerfl"achen).

\begin{columns}
\begin{column}{.5\textwidth}
{\ttfamily {\footnotesize
\color{nounibaredI}\color{nounibaredI}\textbackslash documentclass\color{black}\{article\} \\
\color{nounibaredI}\color{nounibaredI}\textbackslash usepackage\color{black}\{graphicx\} \\
\color{nounibaredI}\color{nounibaredI}\textbackslash usepackage\color{black}\{float\} \\
\color{nounibaredI}\color{unibablueI}\textbackslash\color{unibablueI}begin\color{black}\color{black}\{document\} \\
\color{nounibaredI}\color{unibablueI}\textbackslash\color{unibablueI}begin\color{black}\color{black}\{figure\}\color{nounibagreenI}[H]\color{black} \\
\color{nounibaredI}\color{unibablueI}\textbackslash\color{unibablueI}begin\color{black}\color{black}\{center\} \\
	\color{nounibaredI}\color{nounibaredI}\textbackslash includegraphics\color{black}\color{nounibagreenI}[width=70mm]\color{black}\{path/tux.png\} \\
\color{nounibaredI}\color{nounibaredI}\textbackslash caption\color{black}\{Tiny Tux in Float\} \\
\color{nounibaredI}\color{nounibaredI}\textbackslash label\color{black}\{img:tinytux-float\} \\
\color{nounibaredI}\color{unibablueI}\textbackslash\color{unibablueI}end\color{black}\color{black}\{center\} \\
\color{nounibaredI}\color{unibablueI}\textbackslash\color{unibablueI}end\color{black}\color{black}\{figure\} \\
\color{nounibaredI}\color{unibablueI}\textbackslash\color{unibablueI}end\color{black}\color{black}\{document\} \\
}}
\end{column}

\begin{column}{.5\textwidth} 
\begin{figure}
\begin{center}
    \includegraphics[width=35mm]{image/tux.png}
\caption{Der kleine Tux jetzt in Float}
\label{img:kleinertux_float}
\end{center}
\end{figure}
\end{column}
\end{columns}

\end{frame}
